\documentclass[11pt,a4paper]{article}
\usepackage{fontspec}
\usepackage[margin=2.4cm]{geometry}
\usepackage{graphicx}
\usepackage{xcolor}
\usepackage{booktabs}
\usepackage{titlesec}
\usepackage{parskip}
\usepackage{enumitem}
\usepackage[hidelinks]{hyperref}
\usepackage{fancyhdr}

\setmainfont{DejaVu Serif}[Scale=0.92]
\setsansfont{DejaVu Sans}[Scale=0.92]
\setmonofont{DejaVu Sans Mono}[Scale=0.84]

\definecolor{azul}{HTML}{2a78d6}
\definecolor{naranja}{HTML}{eb6834}
\definecolor{tinta2}{HTML}{52514e}
\definecolor{grisc}{HTML}{f0efec}

\titleformat{\section}{\sffamily\Large\bfseries\color{azul}}{\thesection.}{0.5em}{}
\titleformat{\subsection}{\sffamily\large\bfseries}{}{0em}{}
\titlespacing{\section}{0pt}{18pt}{6pt}

\pagestyle{fancy}\fancyhf{}
\fancyfoot[C]{\sffamily\small\color{tinta2}\thepage}
\fancyhead[L]{\sffamily\footnotesize\color{tinta2}Memoria persistente en modelos de lenguaje}
\fancyhead[R]{\sffamily\footnotesize\color{tinta2}9 de septiembre de 2026}
\renewcommand{\headrulewidth}{0.3pt}

\newcommand{\destacado}[1]{%
  \par\medskip\noindent\colorbox{grisc}{\parbox{\dimexpr\linewidth-2\fboxsep}{%
  \vspace{2pt}#1\vspace{2pt}}}\par\medskip}

\newcommand{\fig}[2]{\begin{center}\includegraphics[width=#2\linewidth]{figuras/#1}\end{center}}

\begin{document}

\begin{center}
{\sffamily\bfseries\LARGE Qué hace cada modelo, explicado desde cero}\\[4pt]
{\sffamily\large\color{tinta2} Micro LM, Micro LM de frontera, TinyLlama y Qwen}\\[10pt]
{\sffamily\small\color{tinta2} Maximiliano Rodrigo Speranza · 9 de septiembre de 2026}
\end{center}

\vspace{6pt}
\hrule
\vspace{14pt}

\section{El problema, en una frase}

Un modelo de lenguaje como los que usa todo el mundo tiene un defecto raro. \textbf{Se olvida de lo
que le dijiste.} Si hoy le contás algo sobre vos, mañana no lo sabe. Y tiene un segundo defecto que
es peor, porque cuando no sabe algo \textbf{no se calla, inventa}, y lo inventa con la misma
seguridad con la que dice las cosas ciertas.

Los dos defectos parecen distintos y son el mismo. El modelo no tiene un lugar donde guardar lo que
le pasó, así que cuando le preguntás algo que debería estar guardado, contesta con lo que más se le
parece de todo lo que leyó alguna vez.

\destacado{El objetivo de esta línea de investigación es una sola cosa. \textbf{Que un modelo de
lenguaje no olvide nunca lo que le dijiste}, y que cuando no lo sepa, lo diga.}

Para resolver eso hace falta entender cómo un modelo guarda y busca información adentro suyo. Y
estudiar eso directamente en un modelo gigante es casi imposible, porque no se le puede apagar lo
que ya sabe. Si le preguntás dónde vive Ana y contesta bien, no hay forma de saber si leyó lo que
vos guardaste o si se acordó de alguna Ana que leyó en internet.

Por eso hay cuatro modelos en juego, y cada uno cumple una función distinta.

\section{Los cuatro modelos}

\fig{01_arquitecturas.pdf}{1.0}

Todos los modelos de lenguaje están hechos de bloques apilados. El texto entra por arriba, pasa por
todos los bloques uno tras otro, y sale la respuesta. Lo que diferencia a un modelo de otro es
cuántos bloques tiene y \textbf{cuánto puede mirar cada bloque}.

Hay dos clases de bloque y la diferencia es toda la historia.

\begin{itemize}[leftmargin=1.4em,itemsep=4pt]
\item \textbf{Atención completa.} El bloque puede mirar todo lo que vino antes. Es potente y es caro,
porque el costo crece con el cuadrado de lo que hay que mirar.
\item \textbf{Regla delta.} El bloque mantiene un resumen que va actualizando, y mira una ventanita.
Es mucho más barato, y es lo que permite contextos larguísimos.
\end{itemize}

\subsection{Micro LM}

Es un modelo entrenado desde cero, de \textbf{3,5 megabytes}. Para dar una idea, una foto del
teléfono pesa más que eso. Sabe 242 palabras en total y contesta de a una palabra por vez.

Suena a juguete y no lo es. \textbf{Es un instrumento de medición.} Como se entrena desde cero y no
sabe nada del mundo, se le puede vaciar la memoria y comprobar que efectivamente contesta cero.
Ese control, que en un modelo grande es imposible, es lo que vuelve interpretable todo lo demás.

La analogía honesta es que el Micro LM no es un auto barato. Es un túnel de viento. Nadie maneja un
túnel de viento hasta la esquina, y sin embargo los aviones se diseñan ahí.

\subsection{Micro LM de frontera}

Es el mismo modelo, con un cambio quirúrgico. En el Micro LM, cuando el modelo va a buscar algo en
su archivo, la pregunta que formula para buscar \textbf{sale de una ventanita} y no ve la pregunta
entera. En el Micro LM de frontera esa consulta pasa a mirar todo.

Ese cambio se hizo para contestar una pregunta concreta, y la respuesta fue mejor que la esperada.

\fig{02_cobertura.pdf}{0.82}

Con una ventana demasiado corta el modelo acierta 0,6430, porque hay una pieza de la pregunta que le
queda literalmente afuera del campo visual. No es que la ignore, no la puede ver. Con una ventana
apenas más larga acierta 0,9977. Y con atención completa, que puede mirar absolutamente todo,
acierta \textbf{exactamente lo mismo}, 0,9977.

\destacado{Ese empate exacto es el resultado. \textbf{Lo que decide no es cuánto mira sino SI llega a
mirar la pieza que importa.} Es una condición de sí o no, no una escala. Mirar de más no compra
nada.}

\subsection{TinyLlama}

Es un modelo real, de verdad, con 1.100 millones de parámetros, entrenado por otra gente sobre
textos reales. Sirve para responder la pregunta que el Micro LM no puede responder solo, que es si
todo lo que se descubrió en el laboratorio también pasa afuera.

Acá el modelo se \textbf{congela entero}. No se le toca ni un parámetro. Lo único que se entrena es
un archivo externo que se le engancha, y que pesa el \textbf{0,192 por ciento} del modelo.

\subsection{Qwen}

Son modelos abiertos y actuales. Entran por una razón que apareció hoy mismo y que cambia el
encuadre de todo el programa.

\section{La frontera se mudó a nuestra arquitectura}

Durante años los modelos grandes usaron atención completa en todas las capas. Eso cambió. Los
modelos abiertos de frontera de 2026, que son Qwen3.5, Qwen3-Next, Kimi Linear y Kimi K3,
\textbf{abandonaron la atención completa} y se pasaron a una mezcla con mayoría de capas de regla
delta.

\fig{07_qwen_capas.pdf}{1.0}

Esto no sale de un artículo de divulgación. Sale del archivo de configuración del modelo, que
cualquiera puede bajar y leer.

\begin{center}\small\sffamily
\begin{tabular}{ll}
\toprule
\texttt{num\_hidden\_layers} & 24 \\
\texttt{full\_attention\_interval} & 4 \\
\texttt{layer\_types} & 3 de regla delta y 1 de atención, repetido 6 veces \\
\texttt{linear\_conv\_kernel\_dim} & \textbf{4}, que es el tamaño de la ventanita \\
\bottomrule
\end{tabular}
\end{center}

\destacado{La consecuencia importa. \textbf{El Micro LM dejó de ser una arquitectura vieja y pasó a
ser una miniatura de lo que la frontera eligió.} Y la ley del empate exacto explica por qué el
patrón de tres a uno funciona, que es porque hace falta que alguna capa alcance a mirar la pieza que
importa, y poner más no compra nada.}

\section{Lo que se midió hoy sobre un modelo real}

\subsection{La tarea, que es la que más importa de todas}

Se le dice un dato. Después se lo corrige. Después se le pregunta.

\begin{center}\sffamily\small
\begin{tabular}{ll}
\toprule
Primero se le dice & \texttt{Ana vive en Córdoba} \\
Después se lo corrige & \texttt{Ana vive en Salta} \\
Y se le pregunta & \texttt{¿Dónde vive Ana?} \\
\midrule
La respuesta correcta es & \texttt{Salta}, porque es la última \\
\bottomrule
\end{tabular}
\end{center}

Parece trivial y no lo es. Es exactamente lo que un asistente tiene que hacer bien para servir de
algo, porque las cosas cambian y la versión vieja de un dato es peor que no tener el dato. Y es
justamente lo que los modelos actuales hacen mal.

Esta tarea estaba resuelta en el Micro LM desde hace semanas. \textbf{Nunca se había probado sobre
un modelo real.}

\subsection{El resultado}

\fig{03_versiones.pdf}{0.92}

El modelo arranca sin idea, cruza la línea del azar alrededor de los 2.000 pasos, y termina
acertando la versión vigente \textbf{0,9366} de las veces. El piso es 0,5000, que es lo que daría
tirar una moneda entre las dos versiones. La diferencia son 24,8 desviaciones estándar, o sea que la
probabilidad de que esto sea casualidad es indistinguible de cero.

\subsection{Y la prueba de que no está haciendo trampa}

Un número solo no prueba nada. Podría estar adivinando por algún atajo. Así que sobre el modelo ya
entrenado se le rompe el archivo de tres maneras distintas.

\fig{04_controles.pdf}{0.82}

\begin{itemize}[leftmargin=1.4em,itemsep=3pt]
\item Con el archivo en su lugar acierta todo.
\item Si se le dan los mismos datos \textbf{mezclados}, para que cada pregunta reciba el archivo de
otra, cae a cero exacto.
\item Si se le vacía el archivo, cero exacto.
\item Si se le apaga la lectura, cero exacto.
\end{itemize}

\destacado{Cero exacto en los tres. \textbf{La respuesta sale del archivo y de ningún otro lado.} No
se la sabía de antes, no la adivina, no hay atajo.}

Y hay un detalle técnico que es el que vuelve válido todo esto. El archivo guarda, junto a cada dato,
una marca de cuándo se dijo. En la versión anterior del experimento esa marca coincidía con el lugar
que el dato ocupaba en la lista, así que era imposible saber si el modelo leía la marca o
simplemente contaba lugares. Eso se arregló esta mañana. Ahora las marcas van desordenadas respecto
del lugar, y el modelo igual acierta. \textbf{Está leyendo la marca de tiempo.}

\section{La lección que salió cara}

\fig{05_presupuesto.pdf}{0.66}

Los dos números de este gráfico son el mismo modelo, la misma tarea y el mismo código. Lo único que
cambia es cuánto tiempo se lo dejó estudiar.

Con 2.000 pasos el resultado era 0,5595, indistinguible de tirar una moneda. Sobre ese número se
construyeron durante la mañana \textbf{tres explicaciones equivocadas seguidas} de por qué la tarea
no se podía aprender. Se llegó a describir un mecanismo detallado y convincente de por qué el modelo
se quedaba mudo, y a diseñar un remedio para curarlo. Las tres explicaciones y el remedio se cayeron
al correr la versión aburrida del experimento, que era simplemente dejarlo estudiar diez veces más.

\destacado{Antes de explicar por qué algo no se aprende, correrlo con el tiempo que usa la tarea
equivalente que ya funciona. \textbf{La hipótesis aburrida va primera.}}

\section{Y después volvió a pasar dos veces}

Con la tarea funcionando, la pregunta siguiente era qué ocurre cuando el mundo se agranda. Se pasó
de ocho nombres distintos a sesenta, y de cuatro datos guardados a ocho. A 2.000 pasos las dos cosas
parecían romper el modelo por completo, y durante un rato eso pasó por hallazgo.

\fig{08_reloj.pdf}{0.86}

No rompían nada. Las tres tareas llegan al mismo techo. Lo único que cambia es cuánto tardan.

\begin{center}\small\sffamily
\begin{tabular}{lccc}
\toprule
tarea & 2.000 pasos & 20.000 pasos & distancia al azar \\
\midrule
contestar la versión vigente & 0,5595 & \textbf{0,9366} & 24,8 $\sigma$ \\
60 nombres, archivo de 8 & 0,1905 & \textbf{0,9900} & 74,2 $\sigma$ \\
60 nombres, archivo de 4 & 0,2143 & \textbf{0,9838} & 73,6 $\sigma$ \\
\bottomrule
\end{tabular}
\end{center}

En los tres casos los controles dan cero exacto, igual que antes. Si se mezcla el archivo, si se lo
vacía o si se apaga la lectura, el modelo no contesta nada.

\destacado{\textbf{Agrandar el problema no baja el techo, corre el reloj.} Ocho nombres llegan al
techo cerca del paso 1.650. Sesenta nombres con un archivo del doble de grande llegan cerca del
8.000. El mismo techo, más tiempo.}

Vale la pena decir esto sin adornos. En un solo día se propusieron cuatro explicaciones distintas de
un fenómeno que no existía, y las cuatro sonaban razonables. La única que resultó cierta era la que
no explicaba nada, que era que faltaba tiempo. Es un recordatorio incómodo de lo fácil que es
construir una teoría encima de una medición corta.

\section{Por qué los Qwen y cuál conviene}

Se midieron los tres candidatos, cargándolos y probándolos, no leyendo sus fichas.

\begin{center}\small\sffamily
\begin{tabular}{lccc}
\toprule
 & TinyLlama & Qwen3-0.6B & Qwen3.5-0.8B \\
\midrule
capas & 22 & 28 & 24 \\
arquitectura & densa & densa & \textbf{híbrida} \\
ventanita corta & no & no & \textbf{sí, de 4} \\
se le puede enganchar el archivo & sí & sí & sí \\
piezas del banco que le sirven & \textbf{161 de 161} & \textbf{161 de 161} & \textbf{161 de 161} \\
licencia & abierta & abierta & abierta \\
\bottomrule
\end{tabular}
\end{center}

El dato que no estaba garantizado y que abarata todo es el de la anteúltima fila. \textbf{El mismo
banco de pruebas corre en los tres modelos sin cambiar una línea de código.} Ya se verificó
corriéndolo en los tres.

Y eso deja armado un experimento que hasta hoy no existía. Qwen3.5 y Qwen3-0.6B son de la misma
familia, del mismo tamaño y con la misma licencia, pero uno es híbrido y el otro es denso. Comparar
esos dos aísla el efecto de la arquitectura sin ninguna otra diferencia de por medio.

\destacado{La decisión es no cambiar de vehículo todavía. TinyLlama tiene el montaje validado y lo
que falta es terminar de medir sobre él. Cuando eso cierre, el vehículo pasa a ser Qwen3.5 con
Qwen3-0.6B de control, porque es el único que tiene adentro el mecanismo del que hablan todas estas
leyes.}

\section{Qué falta}

\begin{enumerate}[leftmargin=1.6em,itemsep=5pt]
\item \textbf{Sacarle la marca de tiempo y ver que se rompe.} Hoy está probado que el modelo resuelve
las versiones. Falta probar que las resuelve \textbf{gracias a la marca}, y para eso hay que sacarla
y ver caer el resultado a la mitad. Es el experimento con más valor por peso de todos.
\item \textbf{Repetirlo con otras dos semillas al azar.} Un resultado con una sola semilla no se
publica.
\item \textbf{Medir la abstención.} Todo lo de hoy mide si contesta bien. Falta medir si sabe callarse
cuando no sabe, que es la otra mitad del objetivo.
\item \textbf{Correr el banco completo, con todo junto y con el tiempo correcto.} Hoy cada pieza
se midió por separado y todas llegan al techo. Falta la corrida con las cuatro cosas a la vez.
\end{enumerate}

\vspace{10pt}
\hrule
\vspace{8pt}
{\footnotesize\color{tinta2}\sffamily
Todos los números de este documento salen de corridas propias, con el código y los datos guardados
en el repositorio del proyecto. Cada resultado se acompaña de sus controles. Las medias se calculan
sobre todas las mediciones de la segunda mitad de cada corrida, nunca sobre la última, y se informa
la distancia al piso en desviaciones estándar.}

\end{document}
